Esta tesis estudia el problema de construir un \zkBridge para \Cardano. En términos generales, un \Bridge busca transferir valor o información entre dos cadenas sin exigir a sus usuarios confianza plena en un intermediario centralizado. En nuestro caso, el objetivo es más específico: diseñar una arquitectura que permita tomar un evento ocurrido en una cadena de origen, autenticarlo criptográficamente y habilitar una acción correspondiente en una cadena destino, todo ello respetando las restricciones reales del entorno de ejecución de \Cardano.

El desafío no es solamente criptográfico. Un \zkBridge práctico debe combinar, al mismo tiempo, un mecanismo de autenticación del estado remoto, una representación compacta de la evidencia relevante, una estrategia de verificación \Onchain compatible con límites estrictos de cómputo y tamaño, y un modelo de estado que preserve invariantes como la atomicidad y la prevención de \textit{replay}. En \Cardano, estas exigencias quedan además condicionadas por el modelo \eUTxO, por los presupuestos de \ExUnits y por la forma en que se expresan e interconectan los validadores.

La propuesta desarrollada en este trabajo parte de una decisión central: no intentar verificar directamente dentro de un circuito \ZK todo el consenso nativo de \Cardano, sino apoyarse en \Mithril como mecanismo de certificación compacta del estado relevante y en \Groth como esquema de verificación sucinta \Onchain. Sobre esa base, la tesis organiza el \zkBridge en una capa de actualización, que mantiene en la cadena destino un estado autenticado y reutilizable, y una capa de aplicación, que consume esa autenticación para autorizar operaciones concretas del protocolo.

Con ese marco, el recorrido del manuscrito avanza desde la formulación del problema hacia las restricciones del ecosistema \Cardano, las decisiones de diseño que esas restricciones fuerzan y, finalmente, una implementación por etapas del camino principal \(\LockTx \rightarrow \MintTx\). Los fundamentos criptográficos y de aritmetización necesarios para sostener esa discusión se desarrollan aparte, de modo que el cuerpo principal pueda concentrarse en la arquitectura del \zkBridge y en sus compromisos de diseño. El lector que necesite revisar esas herramientas puede consultar los apéndices finales dedicados a primitivas criptográficas y construcción de pruebas \SNARK.
